\noindent\textbf{能量项的具体离散形式。}
\label{app:energies}
这些能量项的推导见离散弹性杆模型~\cite{spillmann_corde_2007}。下面简要介绍它们。

离散弯曲和扭转势能计算如下：
\begin{equation}
 V_{i}\left(\mathbf{q}\right) = \frac{1}{2}l_i\sum_{k=1}^3 \mathbf{K}_{kk}\left(\frac{2}{||\mathbf{q}_i||^2} (\mathbf{B}_k(\q_i+\q_{i+1}))^{\top}\frac{\mathbf{q}_{i+1} -  \mathbf{q}_{i}}{l_i+l_{i+1}}- \overline{\epsilon}_k\right)^2,
 \label{eq:V_q}
\end{equation}
其中对角刚度矩阵$\mathbf K$的元素为$\mathbf K_{11} = \mathbf{K}_{22} =  E\pi r^2/4$和$\mathbf K_{33} = G\pi r^2/2$，其中$E$和$G$分别为编码弯曲和扭转抗性的刚度模量。$\overline{\epsilon}_k$是杆基于其静止形状几何的自然弯曲和扭转。$\mathbf B_k \in \mathcal R^{4\times 4}$是常数斜对称矩阵~\cite{spillmann_corde_2007}。

动能由平移动能$T_p(\p)$和转动动能$T_q(\q)$组成。
$T_p\left(\mathbf p \right)$表示隐式欧拉时间离散化下的增量平移动能：
\begin{equation}
  T_p\left(\mathbf p \right)=\frac{1}{2}\left\|\mathbf{p}-\mathbf{p}_{t}-\dot{\mathbf{p}_{t}}\Delta t\right\|_{\mathbf{M} / \Delta t^{2}}^{2},
\end{equation}
其中$\mathbf{M}$是集中质量矩阵，$\Delta t$是时间步长，$\mathbf{p}_{t}$和$\dot{\mathbf{p}}_{t}$分别是上一时间步的顶点位置和速度。
对于截面半径为$r$、材料密度为$\rho$的杆，线段$\mathbf{e}_i$的离散转动动能计算如下：
\begin{equation}
 T_{q,i}\left(\mathbf{q}\right)= \frac{1}{8}l_i\sum_{k=1}^3 \mathbf{J}_{kk}\left( (\mathbf{B}_k(\q_i+\q_{i+1}))^{\top}(\dot{\mathbf{q}}_i+  \dot{\mathbf{q}}_{i+1} ) \right)^2,
 \label{eq:T_q}
\end{equation}
其中对角惯性张量的三个元素分别为$\mathbf{J}_{11}=\mathbf{J}_{22}=\rho\pi r^2/4$和$\mathbf{J}_{33}=\rho\pi r^2/2$。
\iffalse
\begin{equation}
  \setlength{\arraycolsep}{1.2pt}
  \mathbf B_1= \begin{bmatrix}
    0 & 0 & 0 & 1 \\
    0 & 0 & 1 & 0 \\
    0 & -1 & 0 & 0\\
    -1 & 0 & 0 & 0
    \end{bmatrix}, \mathbf B_2= \begin{bmatrix}
    0 & 0 & -1 & 0 \\
    0 & 0 & 0 & 1 \\
    1 & 0 & 0 & 0\\
    0 & -1 & 0 & 0
    \end{bmatrix}, \mathbf B_3= \begin{bmatrix}
    0 & 1 & 0 & 0 \\
    -1 & 0 & 0 & 1 \\
    0 & 0 & 0 & 1\\
    0 & 0 & -1 & 0
    \end{bmatrix} \nonumber
\end{equation}
\fi